% =====================================================================
%  Annexe technique : comprendre les points délicats (tests, cache, auth...)
%  Compilation : pdflatex annexe_technique.tex  (ou Overleaf, ou Tectonic)
% =====================================================================
\documentclass[11pt,a4paper]{article}

\usepackage[utf8]{inputenc}
\usepackage[T1]{fontenc}
\usepackage[french]{babel}
\usepackage{lmodern}
\usepackage{amsmath}
\usepackage{textcomp}
\usepackage[a4paper,margin=2.5cm]{geometry}
\usepackage{xcolor}
\usepackage{booktabs}
\usepackage{enumitem}
\usepackage{listings}
\usepackage[hidelinks]{hyperref}

\definecolor{accent}{RGB}{30,90,150}
\newcommand{\code}[1]{\texttt{#1}}

\lstset{
  basicstyle=\ttfamily\footnotesize,
  backgroundcolor=\color{gray!8},
  frame=single, rulecolor=\color{gray!40},
  breaklines=true, columns=fullflexible,
  showstringspaces=false, xleftmargin=10pt,
}

\setlength{\emergencystretch}{3em}
\sloppy

\title{\vspace{-1.5cm}\textbf{Annexe technique : les points délicats expliqués}\\[2pt]
\large Plateforme de backtesting PRO3600 --- aide à la soutenance}
\author{}
\date{\today}

\begin{document}
\maketitle
\vspace{-0.8cm}

\noindent\textit{Ce document explique, en langage simple, les parties du projet
les plus susceptibles de soulever des questions. Objectif : pouvoir les défendre
sans être pris au dépourvu. La partie sur les tests Vitest est volontairement
détaillée.}

% =====================================================================
\section{Les tests automatisés : l'idée générale}
% =====================================================================
Un \textbf{test automatisé} est un petit programme qui vérifie qu'une partie de
notre code fait bien ce qu'on attend. Au lieu de cliquer à la main dans
l'application à chaque modification, on lance une commande et l'ordinateur
vérifie tout en quelques secondes. S'il affiche « \emph{vert} », tout va bien ;
s'il affiche « \emph{rouge} », quelque chose a cassé.

On distingue deux types :
\begin{itemize}[leftmargin=1.4em,nosep]
  \item \textbf{test unitaire} : on teste \emph{une} fonction isolée
        (ex.\ « ce calcul de rendement donne-t-il le bon nombre~? ») ;
  \item \textbf{test d'intégration} : on teste plusieurs morceaux ensemble
        (ex.\ « quand on appelle l'API pour lancer un backtest, renvoie-t-elle le
        bon résultat~? »).
\end{itemize}
Dans le projet : \textbf{pytest} teste le backend (Python), \textbf{Vitest} teste
le frontend (Angular/TypeScript).

% =====================================================================
\section{Les tests Vitest (frontend) --- en détail}
% =====================================================================
\textbf{Vitest} est l'outil qui exécute nos tests côté frontend. C'est un
« lanceur de tests » : on lui écrit des petites fonctions de vérification, il les
exécute toutes et nous dit combien passent. Dans notre projet Angular, on le
lance avec \code{ng test} (ou \code{npm test}).

\subsection{La structure d'un test (3 mots-clés)}
Un test Vitest s'écrit toujours avec trois éléments :
\begin{itemize}[leftmargin=1.4em,nosep]
  \item \code{describe(...)} : un \textbf{groupe} de tests (ex.\ tous les tests
        d'un service) ;
  \item \code{it(...)} : \textbf{un} test précis, avec une phrase qui décrit ce
        qu'il vérifie ;
  \item \code{expect(...).toBe(...)} : la \textbf{vérification} (« je m'attends à
        ce que ce résultat soit égal à... »).
\end{itemize}

Exemple réel (simplifié) tiré du service d'authentification :
\begin{lstlisting}
describe('AuthService', () => {              // groupe
  it('est authentifie si un user est stocke', () => {   // un test
    localStorage.setItem('auth_user', '{"id":1}');
    const auth = new AuthService();
    // isAuthenticated() doit valoir true puisqu'un user est present
    expect(/* resultat */ true).toBe(true);  // verification
  });
});
\end{lstlisting}
On lit cela comme une phrase : « \emph{dans le groupe AuthService, le test « est
authentifié si un user est stocké » : on met un user dans le stockage, et on
s'attend à ce que isAuthenticated renvoie vrai} ».

\subsection{TestBed : un mini-Angular pour les tests}
Un composant Angular (une page, par exemple) a besoin de ses \emph{dépendances}
pour fonctionner : le routeur (pour naviguer), le client HTTP (pour parler au
serveur), etc. Dans un vrai navigateur, Angular les fournit automatiquement. Dans
un test, il faut les fournir nous-mêmes : c'est le rôle de \textbf{TestBed}, qui
monte un « mini-Angular » de laboratoire. D'où des lignes comme :
\begin{lstlisting}
TestBed.configureTestingModule({
  providers: [provideRouter([]), provideHttpClient(), provideHttpClientTesting()],
});
const fixture = TestBed.createComponent(Backtests);  // cree la page en labo
\end{lstlisting}
Traduction : « prépare un mini-Angular avec un faux routeur et un faux client
HTTP, puis crée la page Backtests dedans ». C'est pour cela que, sans ces
\emph{providers}, certains tests échouaient avec « \emph{No provider for
ActivatedRoute} » : la page réclamait le routeur, qu'on avait oublié de fournir.

\subsection{HttpTestingController : tester sans vrai serveur}
Nos services appellent l'API (ex.\ \code{POST /api/backtest}). On ne veut pas
qu'un test dépende d'un vrai serveur allumé (ce serait lent et fragile). On
utilise donc \textbf{HttpTestingController}, qui \emph{intercepte} les appels HTTP
et nous laisse fournir une fausse réponse :
\begin{lstlisting}
service.runBacktest(requete).subscribe();
const appel = httpMock.expectOne('.../api/backtest'); // on capture l'appel
expect(appel.request.method).toBe('POST');            // bon verbe HTTP ?
appel.flush({});                                       // fausse reponse renvoyee
\end{lstlisting}
Traduction : « j'appelle runBacktest ; je vérifie qu'un \code{POST} part bien vers
\code{/api/backtest} ; et je réponds avec un objet vide à la place du serveur ».
On teste donc que \emph{notre} code appelle la bonne URL avec la bonne méthode,
sans réseau.

\subsection{Pourquoi on ne dessine pas les graphiques dans les tests}
Les tests tournent dans \textbf{jsdom}, un « faux navigateur » en mémoire (sans
écran réel). Or la librairie de graphiques (\code{lightweight-charts}) a besoin
d'un vrai canvas pour dessiner, ce que jsdom n'a pas. Nos tests évitent donc de
déclencher l'affichage des graphiques (par exemple en ne fournissant pas de
résultat à afficher). On teste la \emph{logique}, pas le rendu visuel du
graphique.

\subsection{Ce que vérifient concrètement nos tests frontend}
\begin{itemize}[leftmargin=1.4em,nosep]
  \item \textbf{AuthService} : un utilisateur stocké rend l'application
        « connectée » ; le mode démo aussi ; la déconnexion vide le stockage.
  \item \textbf{BacktestHistoryService} : l'historique local garde au maximum les
        5 derniers backtests (on en ajoute 7, on vérifie qu'il n'en reste que 5).
  \item \textbf{DataService} : chaque méthode appelle la bonne URL avec la bonne
        méthode HTTP (GET, POST, DELETE...).
  \item \textbf{Garde de route} (\code{authGuard}) : si on n'est pas connecté, on
        est bien redirigé vers \code{/login}.
  \item \textbf{Pages} (backtests, results, dashboard...) : elles se créent sans
        erreur et leurs petites fonctions (formatage de pourcentage, choix des
        paramètres par défaut d'une stratégie, explication d'une stratégie...)
        renvoient les bonnes valeurs.
\end{itemize}

\subsection{Le « scaffolding » qu'il a fallu réparer}
Au départ, \code{ng test} refusait de compiler les fichiers de tests
(\code{*.spec.ts}) : il manquait un fichier de configuration
(\code{tsconfig.spec.json}) indiquant à TypeScript d'inclure les tests, et les
types de Vitest. Une fois ces deux éléments ajoutés, les tests compilent et
s'exécutent. \emph{À retenir} : si on demande « pourquoi avoir touché à la config
de test~? », la réponse est « les tests n'étaient pas branchés, on a ajouté le
\code{tsconfig.spec.json} manquant ».

\subsection{Comment les lancer}
\begin{lstlisting}
cd frontend
npm test            # equivaut a : ng test --no-watch
\end{lstlisting}
Le résultat ressemble à \code{Test Files 11 passed (11) / Tests 41 passed (41)} :
cela veut dire que les 41 vérifications réparties dans 11 fichiers passent toutes.

% =====================================================================
\section{Les tests pytest (backend) --- en bref}
% =====================================================================
Même principe, côté Python, avec \textbf{pytest}. Deux idées valent la peine
d'être comprises :
\begin{itemize}[leftmargin=1.4em,nosep]
  \item \textbf{Scénario « calculé à la main »} : pour le moteur de backtest, on
        fournit une mini-série de prix dont on a calculé soi-même, sur papier, le
        rendement, le ratio de Sharpe, etc. Le test vérifie que le code retrouve
        exactement ces nombres. C'est la preuve que les formules sont justes.
  \item \textbf{« Mock » de Yahoo Finance} : on remplace l'appel réseau qui
        télécharge les vrais cours par une fausse fonction qui renvoie des données
        fabriquées. Ainsi les tests sont rapides, reproductibles et ne dépendent
        pas d'Internet ni de la base de production (on utilise une petite base
        SQLite temporaire).
\end{itemize}
Lancement : \code{cd backend} puis \code{pytest} (57 tests).

% =====================================================================
\section{Autres points « chauds » expliqués}
% =====================================================================
\begin{description}[leftmargin=1.2em,style=nextline]
  \item[Éviter de « tricher » (look-ahead).] Dans un backtest, on ne doit jamais
        utiliser une information qu'on ne connaîtrait pas encore en réalité. Notre
        moteur décale donc le signal d'une journée (\code{signal.shift(1)}) : la
        décision prise aujourd'hui s'applique au rendement de demain, pas
        d'aujourd'hui.
  \item[Le cache des cours (24\,h).] Pour ne pas redemander sans cesse les mêmes
        données à Yahoo Finance, on les stocke en base (table
        \code{market\_data\_cache}) pendant 24~heures. Si on relance le même
        ticker/période dans la journée, on lit le cache au lieu d'appeler le
        réseau --- c'est plus rapide et plus fiable.
  \item[Connexion Google (OAuth) et mode démo.] OAuth est le mécanisme standard
        « se connecter avec Google » : Google authentifie l'utilisateur et nous
        renvoie son identité, sans qu'on gère de mot de passe. Le \emph{mode démo}
        permet d'utiliser l'application sans compte (les données restent alors
        dans le navigateur).
  \item[Le mode avancé (import de trades).] En plus des stratégies intégrées,
        l'utilisateur peut envoyer ses propres transactions (fichier JSON) via
        \code{/api/backtest/upload}. Le serveur les valide puis calcule des
        métriques financières plus poussées : ratio de Sharpe, de \textbf{Sortino}
        (variante qui ne pénalise que les pertes), \textbf{profit factor} (gains
        totaux / pertes totales), espérance par trade, etc.
  \item[Angular « standalone » et « signals ».] Notre application n'utilise pas
        l'ancien système de modules d'Angular (\emph{standalone}) et gère sa
        réactivité avec des \emph{signals} (des variables qui préviennent
        automatiquement l'interface quand elles changent). C'est la manière
        moderne (Angular~21).
  \item[ORM SQLAlchemy.] Plutôt que d'écrire du SQL à la main, on décrit nos
        tables comme des classes Python (\code{User}, \code{Backtest}...) et la
        librairie SQLAlchemy se charge de dialoguer avec la base. Cela évite les
        erreurs et fonctionne aussi bien avec SQLite (local) que PostgreSQL (prod).
\end{description}

\bigskip
\noindent\textit{En résumé pour la soutenance : nous avons des tests unitaires et
d'intégration, lancés par pytest (backend) et Vitest (frontend), qui passent tous.
Les choix « délicats » (anti-triche du backtest, cache, OAuth, mode avancé) ont
chacun une raison simple et pratique, rappelée ci-dessus.}

\end{document}
